\documentclass[12pt,a4paper]{article}
\usepackage[utf8]{vietnam}
\usepackage{amsmath,amssymb}
\usepackage{geometry}
\geometry{
  top=12mm,
  bottom=12mm,
  left=14mm,
  right=14mm,
  footskip=6mm,
  headheight=14pt
}

\usepackage{tikz}
\usetikzlibrary{arrows.meta,calc,shapes.geometric}
\usepackage[most]{tcolorbox}
\usepackage{fontawesome}
\usepackage{fancyhdr}
\usepackage{enumitem}

\pagestyle{fancy}
\fancyhf{}
\fancyhead[L]{\footnotesize\color{teal!80!black}\faBook\ \textbf{SỔ TAY KIẾN THỨC TOÁN 3}}
\fancyhead[R]{\footnotesize\color{teal!80!black}\textbf{Lớp Toán Thầy Lĩnh -- loptoanthaylinh.vn}}
\fancyfoot[C]{\footnotesize\color{gray}-- Sổ tay kiến thức: Bài 12. Bảng nhân 9, bảng chia 9 --}
\renewcommand{\headrulewidth}{0.4pt}
\renewcommand{\footrulewidth}{0pt}

\setlength{\parindent}{0pt}
\setlength{\parskip}{0pt}

\begin{document}

% --- BANNER BÀI HỌC ---
\begin{tcolorbox}[
  enhanced,
  colback=teal!8!white,
  colframe=teal!70!black,
  arc=3mm, boxrule=1.2pt,
  left=3mm, right=3mm, top=2mm, bottom=2mm,
  halign=center
]
  {\LARGE\bfseries\color{teal!90!black}\faBookmark\ BÀI 12: BẢNG NHÂN 9, BẢNG CHIA 9}\\[1mm]
  {\itshape\footnotesize\color{black!75} Quy tắc ngón tay kì diệu, tổng các chữ số của tích luôn bằng 9 trên 1 trang A4}
\end{tcolorbox}

\vspace{-2mm}

\noindent
\begin{minipage}[t]{0.49\textwidth}\vspace{0pt}
\begin{tcolorbox}[
  enhanced,
  colback=blue!4!white,
  colframe=blue!75!black,
  arc=2mm, boxrule=0.8pt,
  left=2mm, right=2mm, top=1.8mm, bottom=1.8mm,
  title={\footnotesize\bfseries\color{white}\faTimes\ 1. Bảng nhân 9 (Ảo thuật số học)}
]
{\footnotesize
\begin{itemize}\setlength{\itemsep}{0.6pt}\setlength{\leftskip}{-2.5mm}
  \item Dãy tích: $9, 18, 27, 36, 45, 54, 63, 72, 81, 90$.
  \item \textbf{Quy luật tổng bằng 9 tuyệt đỉnh:}
  Trong tích $9 \times k$, \textbf{tổng 2 chữ số luôn bằng 9}:
  $$1+8 = 2+7 = 3+6 = 4+5 = \mathbf{9}$$
  $$5+4 = 6+3 = 7+2 = 8+1 = \mathbf{9}$$
  \item \textbf{Chữ số hàng chục} luôn kém số nhân $k$ đúng $1$ đơn vị:
  $$9 \times \mathbf{7} = \mathbf{6}3\quad (6 = 7 - 1;\ 6 + 3 = 9)$$
  $$9 \times \mathbf{8} = \mathbf{7}2\quad (7 = 8 - 1;\ 7 + 2 = 9)$$
\end{itemize}
}
\end{tcolorbox}
\end{minipage}
\hfill
\begin{minipage}[t]{0.49\textwidth}\vspace{0pt}
\begin{tcolorbox}[
  enhanced,
  colback=teal!4!white,
  colframe=teal!75!black,
  arc=2mm, boxrule=0.8pt,
  left=2mm, right=2mm, top=1.8mm, bottom=1.8mm,
  title={\footnotesize\bfseries\color{white}\faShareAlt\ 2. Bảng chia 9 (Từ nhân suy ra chia)}
]
{\footnotesize
\begin{itemize}\setlength{\itemsep}{0.6pt}\setlength{\leftskip}{-2.5mm}
  \item Lấy tích chia cho 9 thì được thừa số kia:
  \begin{tabular}{ll}
    $9 : 9 = 1$ & $54 : 9 = 6$\\
    $18 : 9 = 2$ & $63 : 9 = 7$\\
    $27 : 9 = 3$ & $72 : 9 = 8$\\
    $36 : 9 = 4$ & $81 : 9 = 9$\\
    $45 : 9 = 5$ & $90 : 9 = 10$
  \end{tabular}
  \item Mẹo tính chia: Với số có 2 chữ số chia cho 9, thương chính bằng \textbf{chữ số hàng chục cộng thêm 1} (VD: $72 : 9 \to 7 + 1 = \mathbf{8}$).
\end{itemize}
}
\end{tcolorbox}
\end{minipage}

\vspace{-2.5mm}
% --- KHỐI 3: SƠ ĐỒ TƯ DUY TIKZ ---
\noindent
\begin{tcolorbox}[
  enhanced,
  colback=orange!4!white,
  colframe=orange!80!black,
  arc=2mm, boxrule=0.8pt,
  left=2mm, right=2mm, top=1.5mm, bottom=1.5mm,
  title={\footnotesize\bfseries\color{white}\faSitemap\ SƠ ĐỒ TƯ DUY: QUY TẮC BÀN TAY KÌ DIỆU \& BẢNG NHÂN 9}
]
\centering
\begin{tikzpicture}[
  box/.style={draw, rounded corners=2pt, align=center, font=\scriptsize\bfseries, inner sep=2.5pt, minimum height=0.75cm},
  arr/.style={-{Stealth[length=2mm]}, thick}
]
  \node[box, fill=purple!15, draw=purple!80!black, text width=3.3cm] (p1) at (0, 0) {BẬP NGÓN THỨ $k$\\Ví dụ $9 \times 4$:\\Gập ngón thứ 4};
  \node[box, fill=blue!12, draw=blue!80!black, text width=3.3cm] (p2) at (4.7, 0) {BÊN TRÁI: HÀNG CHỤC\\Có $3$ ngón\\$\implies \mathbf{30}$};
  \node[box, fill=teal!15, draw=teal!80!black, text width=3.3cm] (p3) at (9.4, 0) {BÊN PHẢI: ĐƠN VỊ\\Có $6$ ngón\\$\implies \mathbf{6} \to 9 \times 4 = \mathbf{36}$};

  \draw[arr, purple!80!black] (p1) -- (p2);
  \draw[arr, blue!80!black] (p2) -- (p3);
\end{tikzpicture}
\end{tcolorbox}

\vspace{-2.5mm}
% --- KHỐI 4: VÍ DỤ & MẸO TRÁNH BẪY ---
\noindent
\begin{minipage}[t]{0.49\textwidth}\vspace{0pt}
\begin{tcolorbox}[
  enhanced,
  colback=green!4!white,
  colframe=green!60!black,
  arc=2mm, boxrule=0.8pt,
  left=2mm, right=2mm, top=1.8mm, bottom=1.8mm,
  title={\footnotesize\bfseries\color{white}\faCheckCircle\ Ví dụ thực hành chuẩn mực}
]
{\footnotesize
\textbf{1. Tính nhẩm bằng quy tắc 9:}\\
$9 \times 6$: Chục là $6 - 1 = 5$; Đơn vị là $9 - 5 = 4 \implies \mathbf{54}$.\\
$9 \times 9$: Chục là $9 - 1 = 8$; Đơn vị là $9 - 8 = 1 \implies \mathbf{81}$.\\
\textbf{2. Tính nhẩm phép chia:}\\
$45 : 9 = 4 + 1 = \mathbf{5}$.\\
$63 : 9 = 6 + 1 = \mathbf{7}$.\\
\textbf{3. Bài toán thực tế:}\\
Chia $72$ viên bi đều cho $9$ bạn. Mỗi bạn nhận được: $72 : 9 = \mathbf{8}$ (viên bi).
}
\end{tcolorbox}
\end{minipage}
\hfill
\begin{minipage}[t]{0.49\textwidth}\vspace{0pt}
\begin{tcolorbox}[
  enhanced,
  colback=red!4!white,
  colframe=red!70!black,
  arc=2mm, boxrule=0.8pt,
  left=2mm, right=2mm, top=1.8mm, bottom=1.8mm,
  title={\footnotesize\bfseries\color{white}\faExclamationTriangle\ Mẹo tránh bẫy trắc nghiệm}
]
{\footnotesize
\begin{itemize}\setlength{\itemsep}{1.5pt}\setlength{\leftskip}{-2.5mm}
  \item \textbf{Cặp số đảo ngược đối xứng:}\\
  $9 \times 2 = 18 \leftrightarrow 9 \times 9 = 81$.\\
  $9 \times 3 = 27 \leftrightarrow 9 \times 8 = 72$.\\
  $9 \times 4 = 36 \leftrightarrow 9 \times 7 = 63$.\\
  $9 \times 5 = 45 \leftrightarrow 9 \times 6 = 54$.\\
  Ghi nhớ các cặp số lộn ngược này sẽ không bao giờ viết sai!
\end{itemize}
}
\end{tcolorbox}
\end{minipage}


\end{document}
